% TOP_PROBLEM: {"id": 30006684, "title": "Hilbert's Tenth Problem over the Rationals", "category_id": 1, "category": {"id": 1, "name": "number_theory", "display_name": "Number Theory", "description": "Properties of integers, prime numbers, Diophantine equations.", "slug": "number-theory", "order_index": 1, "created_at": "2026-07-31T15:26:25.670Z"}, "status": "open"}
% FIRST_POSTED: null
% !TeX program = lualatex
% Compile twice: lualatex open_problems_500.tex
% All references and prose are embedded; no BibTeX database or external inputs.
\documentclass[11pt,a4paper]{article}
\usepackage[margin=24mm,headheight=14pt]{geometry}
\usepackage{fontspec}
\setmainfont{Latin Modern Roman}
\setsansfont{Latin Modern Sans}
\setmonofont{Latin Modern Mono}
\usepackage{amsmath,amssymb,mathtools}
\usepackage{unicode-math}
\setmathfont{Latin Modern Math}
\usepackage{microtype}
\usepackage{enumitem,xurl,needspace}
\usepackage[dvipsnames]{xcolor}
\usepackage{hyperref}
\hypersetup{unicode=true,colorlinks=true,linkcolor=MidnightBlue,urlcolor=MidnightBlue,
 pdftitle={Mathematical Research Notebook},pdfauthor={Source-annotated compilation},bookmarksnumbered=true}
\usepackage{bookmark}
\usepackage{tocloft}
\setlength{\cftsecnumwidth}{3em}
\cftsetpnumwidth{2.5em}
\cftsetrmarg{4em}
\usepackage{titlesec}
\titleformat{\section}{\Large\bfseries\raggedright}{\thesection}{0.7em}{}
\usepackage{fancyhdr}
\pagestyle{fancy}\fancyhf{}
\fancyhead[L]{\small\sffamily Open Problems Math | Research notebook}
\fancyhead[R]{\small 27 September 2026}
\fancyfoot[C]{\thepage}
\setlength{\parindent}{0pt}
\setlength{\parskip}{4pt plus 1pt}
\setlength{\emergencystretch}{3em}
\setcounter{tocdepth}{1}
\setcounter{secnumdepth}{1}
\setlist[itemize]{leftmargin=3em,itemsep=3pt,topsep=4pt,parsep=0pt}
\allowdisplaybreaks
\newcommand{\N}{\mathbb{N}}
\newcommand{\Z}{\mathbb{Z}}
\newcommand{\Q}{\mathbb{Q}}
\newcommand{\R}{\mathbb{R}}
\newcommand{\C}{\mathbb{C}}
\newcommand{\F}{\mathbb{F}}
\newcommand{\A}{\mathbb{A}}
\newcommand{\PP}{\mathbb P}
\newcommand{\E}{\mathbb{E}}
\newcommand{\eps}{\varepsilon}
\newcommand{\dd}{\,\mathrm d}
\newcommand{\sref}[2]{\hyperlink{source:#1:#2}{[S#2]}}
\newcommand{\eref}[2]{\hyperlink{extra:#1:#2}{[E#2]}}
% Turn the next switch off for a shorter reading copy. The quotations preserve
% every exactTarget field, including qualifications omitted from a short summary.
\newif\ifshowcatalogue
\showcataloguetrue
\newcommand{\cataloguescope}[1]{\ifshowcatalogue
 \paragraph{Exact scope in the supplied catalogue.}
 \begin{quote}\small #1\end{quote}\fi}
\usepackage{etoolbox}
\AtBeginEnvironment{itemize}{\small}
% Standalone entry; original section content is delimited for verification.
\begin{document}
\setcounter{section}{0}
%% BEGIN ORIGINAL SECTION CONTAINER
% ================================================================
% problemId: 30006684
% Canonical problem ID: 30006684
% exactTarget SHA-256: f045a471c959f3cb88dfc01e85c12b13f46b915f3d6d6f5135d6858aa366d230
\clearpage
\section*{\texorpdfstring{Hilbert's Tenth Problem over the Rationals}{Hilbert's Tenth Problem over the Rationals}}\label{30006684}
\subsection{Definitions and mathematical statement}
An input is a finite description of $f\in\Z[x_1,\ldots,x_n]$, with no fixed bound on $n$, degree, or coefficient size. Define the decision set
\[
 \mathcal D_{\Q}=\{\langle f\rangle:\exists(q_1,\ldots,q_n)\in\Q^n,
                                  f(q_1,\ldots,q_n)=0\}.
\]
Determine whether $\mathcal D_{\Q}$ is decidable: one algorithm must halt on every input and return its correct membership bit. Enumeration of rational solutions only semidecides yes-instances and is not a solution. The selected assertion concerns computability, with no polynomial-time demand.
\par\smallskip\textit{Formulation and concept references:} \sref{30006684}{1}.
\cataloguescope{There exists a total algorithm which, given any finite encoding of a multivariate polynomial f in Z[x\_1,...,x\_n] with arbitrary finite n and integer coefficients, halts and correctly decides whether there is q in Q\textasciicircum{}n with f(q)=0. Standard computably interconvertible polynomial encodings are equivalent; no uniform running-time complexity bound is part of the assertion.}
\subsection{Short English statement}
Can an algorithm always decide whether a polynomial equation with integer coefficients has a rational solution?
%% BEGIN RESEARCH INSERTION
\subsection{Research attempt}
\paragraph{Current frontier.}
\textit{Research check: 27 September 2026.}
The 2026 Koymans--Pagano survey states undecidability over infinite
finitely generated rings. It does not include $\Q$: a ring generated by
finitely many rationals inverts only finitely many primes. \sref{30006684}{2}.
Garcia-Fritz--Pasten--Vidaux prove undecidability when rational equations
are augmented with height-comparison predicates, while explicitly leaving
the ordinary rational problem open. Those predicates cannot silently be
added to this entry. \sref{30006684}{3}.
The local-global survey describes the positive-existential definability
route and its limitations. \sref{30006684}{1}.

\paragraph{Attack route.}
A negative answer could follow from a Diophantine interpretation of
integer arithmetic. A positive answer could follow from a total
computable bound for the least rational solution. A third idea is to
exhaust localizations with finitely many inverted primes. We investigate
the exact height-bound equivalence and use explicit equations to stress
test the latter two strategies.

\paragraph{Attempt: computable height versus a decision procedure.}
Fix a decidable binary encoding of integer polynomials. For
$q=a/b$ in lowest terms with $b>0$, set
$H(q)=\max(|a|,b)$ and $H(q_1,\ldots,q_n)=\max(1,H(q_1),\ldots,H(q_n))$.
The following two statements are equivalent:
\[
 \mathcal D_{\Q}\text{ is decidable};
 \tag{21.1}
\]
\[
 \begin{gathered}
 \text{there is a total computable }B:\N\to\N\text{ such that}\\
 f\in\mathcal D_{\Q},\ |\langle f\rangle|\leq L
 \ \Longrightarrow\ \exists q:\ f(q)=0,\ H(q)\leq B(L).
 \end{gathered}
 \tag{21.2}
\]
If (21.2) holds, enumerate all reduced fractions with numerator absolute
value and positive denominator at most $B(L)$, and all tuples of the
required length. Exact rational arithmetic decides whether any is a
zero. The list is finite, and (21.2) justifies returning no at its end.

Conversely, given a decider in (21.1), list the finitely many valid
polynomial encodings of length at most $L$. Decide each. For every
positive instance, enumerate rational tuples in increasing height until
one solution is found. Each of these finitely many searches terminates;
take the maximum of the found heights and 1. This algorithm computes
$B(L)$. Neither its running time nor the size of $B(L)$ is controlled.
Thus this is a computability equivalence, exactly at the strength of the
selected target, not an effective numerical estimate.

\paragraph{Attempt: reduction to quartic sums of squares.}
Every polynomial can be evaluated by a finite straight-line computation
using addition and multiplication with integer constants. Give each gate
an auxiliary variable $z_j$. Its consistency equation has degree at most
two: $z_j-z_a-z_b=0$, $z_j-z_az_b=0$, or $z_j-c=0$.
Include the equation setting the final value to zero. If the resulting
equations are $g_1=\cdots=g_r=0$, replace them by
\[
 F(x,z)=\sum_{j=1}^r g_j(x,z)^2=0.
 \tag{21.3}
\]
Over $\Q\subset\R$, a sum of squares is zero exactly when every summand
is zero. The auxiliary variables are uniquely determined by the inputs
in gate order. Hence $F$ has a rational zero if and only if the original
polynomial does, and $\deg F\leq4$. This reduction is effective for all
finite inputs; a dense or sparse encoding may change its length but not
its computability. A decider for this quartic sum-of-squares subclass
would therefore decide the original problem. No fixed bound on the
number of variables is asserted.

Now test a proposed small height bound on the especially simple family
\[
 F_r(x_0,\ldots,x_r)=(x_0-2)^2+
              \sum_{j=0}^{r-1}(x_{j+1}-x_j^2)^2.
 \tag{21.4}
\]
Its unique rational zero has $x_j=2^{2^j}$, by nonnegativity followed by
induction. Thus the least multiplicative height is $2^{2^r}$ even though
the degree is four and the sparse description has length $O(r\log(r+2))$.
In particular no bound of the form $2^{L^A}$ with fixed $A$ works for all
these descriptions: $2^r$ eventually exceeds every fixed power of
$r\log(r+2)$. This is not evidence that every computable bound fails;
a doubly exponential computable bound easily accommodates this family.

\paragraph{Attempt: why finite localizations do not finish the search.}
Let $S$ be a finite prime set and $R_S=\Z[p^{-1}:p\in S]$. If
$p\notin S$, the polynomial $px-1$ has the rational zero $1/p$, but no
zero in $R_S$, since every reduced denominator in $R_S$ is supported on
$S$. Therefore a negative answer for any one finite localization cannot
be promoted to a negative answer over $\Q$. Exhausting the primes
recovers a semidecision for rational solutions, but no stopping rule for
negative instances. The fact that $\Q$ is a union of these rings does
not transfer their undecidability to the union or provide its decider.

For the opposing route, an explicit polynomial
$D(t,u_1,\ldots,u_s)$ defining $t\in\Z$ by existential rational
quantification would let us replace each integer variable of an integer
Diophantine equation by its rational counterpart together with $D=0$.
Combining all equations as in (21.3) would yield a reduction of integer
undecidability to this target. The needed existential definition has not
been constructed here; a definition containing universal quantifiers or
height predicates does not meet this requirement. \sref{30006684}{1},
\sref{30006684}{3}.

\paragraph{Stress test.}
The equivalence (21.1)--(21.2) bounds at least one point on each nonempty
variety, not every point. Finite rational sets and varieties with
infinitely many points are handled in exactly the same way. Zero
polynomials and constant nonzero polynomials cause no difficulty.
No compactness argument has converted pointwise eventual discovery into
an effective bound. The quartic reduction uses order only as a proof
that rational squares cannot cancel; it does not add an inequality
oracle to the input language.

\paragraph{Outcome.}
\textbf{Outcome: Conditional reduction.}
The attempt gives an exact reduction to quartic sums of squares and an
exact computable-minimal-height criterion, with a double-exponential
height stress family. These reformulations are not claimed to be new
characterizations. Neither the computable bound nor an existential
integer interpretation is established. The selected assertion of a
total decider remains undecided by this notebook.
%% END RESEARCH INSERTION
\subsection{Sources}
\begin{itemize}\raggedright
\item[\textnormal{[S1]}]\hypertarget{source:30006684:1}{} Sylvy Anscombe, Valentijn Karemaker, Zeynep Kisakürek, Vlerë Mehmeti, Margherita Pagano and Laura Paladino, A survey of local-global methods for Hilbert's Tenth Problem, abstract and Section1.
\url{https://arxiv.org/pdf/2309.14987}.
{\small\textit{Catalogue formulation source.}}
\item[\textnormal{[S2]}]\hypertarget{source:30006684:2}{} Hilbert's tenth problem for finitely generated rings.
\url{https://arxiv.org/html/2602.04468v1}.
{\small\textit{Catalogue research/status reference; not a proof endorsement.}}
\item[\textnormal{[S3]}]\hypertarget{source:30006684:3}{} Effectivity for existence of rational points is undecidable.
\url{https://arxiv.org/abs/2311.01958}.
{\small\textit{Catalogue research/status reference; not a proof endorsement.}}
\end{itemize}

%% END ORIGINAL SECTION CONTAINER
\end{document}
